\section{Implement a Stack using an Array}
\label{sec:sq-stack-array}

\problemstatement{Implement a stack supporting $\textit{push}(x)$,
$\textit{pop}()$, $\textit{top}()$, and $\textit{isEmpty}()$, backed by a
fixed-size array, each operation running in $O(1)$.}

\begin{patternbox}
No pattern recognition is needed here -- this is foundational machinery.
The only idea worth internalizing is that a stack needs just \textbf{one}
extra piece of state beyond the array itself: an index, conventionally
called \textit{top}, pointing at the most recently pushed element.
\end{patternbox}

\subsection*{Understanding the problem carefully}

Since a stack only ever adds or removes from one end (the ``top''), an
array can simulate it perfectly by treating one end (conventionally the
highest used index) as that top, and tracking that index explicitly. No
shifting of other elements is ever required, which is exactly what keeps
every operation $O(1)$.

\begin{ideabox}[The Single Pointer Invariant]
Maintain $\textit{top}$ as the index of the most recently pushed element
($-1$ when empty). $\textit{push}(x)$ writes $x$ at index $\textit{top}+1$
and increments \textit{top}. $\textit{pop}()$ reads $\textit{arr}
[\textit{top}]$ and decrements \textit{top} (the old value is left in
place but is considered logically removed, since future pushes will
overwrite it).
\end{ideabox}

\subsection*{Algorithm}

\begin{enumerate}
  \item \textit{top} $\gets -1$ initially.
  \item $\textit{push}(x)$: increment \textit{top}; set
        $\textit{arr}[\textit{top}] \gets x$.
  \item $\textit{pop}()$: read $\textit{arr}[\textit{top}]$; decrement
        \textit{top}; return the read value.
  \item $\textit{top}()$: return $\textit{arr}[\textit{top}]$.
  \item $\textit{isEmpty}()$: return $(\textit{top} = -1)$.
\end{enumerate}

\subsection*{Dry run}

\dryrun{Push $10, 20, 30$; pop once; push $40$.}

\begin{itemize}
  \item Push $10$: \textit{top}$=0$, $\textit{arr}=[10]$.
  \item Push $20$: \textit{top}$=1$, $\textit{arr}=[10,20]$.
  \item Push $30$: \textit{top}$=2$, $\textit{arr}=[10,20,30]$.
  \item Pop: returns $30$; \textit{top}$=1$.
  \item Push $40$: \textit{top}$=2$, $\textit{arr}=[10,20,40]$ (index $2$
        is overwritten).
\end{itemize}

\subsection*{C++ implementation}

\begin{lstlisting}
#include <bits/stdc++.h>
using namespace std;

class Stack {
    vector<int> arr;
    int topIdx;

public:
    Stack(int capacity) : arr(capacity), topIdx(-1) {}

    void push(int x) {
        arr[++topIdx] = x;
    }

    int pop() {
        return arr[topIdx--];
    }

    int top() {
        return arr[topIdx];
    }

    bool isEmpty() {
        return topIdx == -1;
    }
};

int main() {
    Stack s(10);
    s.push(10); s.push(20); s.push(30);
    cout << "Popped: " << s.pop() << endl;
    s.push(40);
    cout << "Top: " << s.top() << endl;
    return 0;
}
\end{lstlisting}

\begin{complexitybox}
\textbf{Time Complexity.} $O(1)$ for every operation -- only a single
index update and a single array access are involved.
\textbf{Space Complexity.} $O(\textit{capacity})$ for the underlying
array.
\end{complexitybox}

\begin{takeawaybox}
A stack needs exactly one pointer into a contiguous array. Compare this to
Section~\ref{sec:sq-queue-linked-list}, where a queue's two-ended access
pattern makes a plain array awkward (without wraparound) and motivates a
linked-list-based implementation instead.
\end{takeawaybox}
